\noindent
{\Large \textbf{\section{Математическое описание}}} \par
\vspace{-20pt}

{\Large \textbf{\subsection{Конечный автомат}}} \par

Конечным автоматом называется набор
\[
 A = (S, \Sigma, Y, s_0, \delta, \lambda),
\]
где
\begin{itemize}[itemsep=-3pt]
 \item $S$ - конечное множество состояний
 \item $\Sigma$ - конечное множество входных сигналов
 \item $Y$ - конечное множество выходных сигналов
 \item $s_0 \in S$ - начальное состояние
 \item $\delta: S \times \Sigma \to S$ - функция переходов
 \item $\lambda: S \times \Sigma \to Y$ - функция выходов
\end{itemize}

\vspace{-30pt}

{\Large \textbf{\subsection{Клеточный автомат}}} \par

Клеточный автомат - это дискретная модель, состоящая из регулярной решетки клеток, каждая из которых в момент времени $t$ находится в одном из конечных состояний. Состояние клетки на шаге $t+1$ определяется её состоянием и состояниями клеток её окрестности на шаге $t$ по одному и тому же правилу для всех клеток. Обновление всех клеток происходит одновременно.

Клеточный автомат формально определяется как
\[
 A = (L, S, \mathcal{N}, f),
\]
где
\begin{itemize}[itemsep=-3pt]
 \item $L=\{1,\dots,n\}\times\{1,\dots,n\}$ - конечная двумерная решетка размера $n\times n$ с тороидальными граничными условиями
 \item $S=\{0,1\}$ - множество состояний клетки (0 - мертвая, 1 - живая)
 \item $\mathcal{N}$ - окрестность фон Неймана
 \item $f: S^5\to S$ - функция перехода
\end{itemize}
%Начальное состояние - отображение $s^0: L\to S$, задающее конфигурацию поля при $t=0$ (случайную или заданную вручную).

%Каждая клетка является конечным автоматом из п. 2.1, у которого входным сигналом $\Sigma$ служат состояния клеток окрестности, функция переходов $\delta$ совпадает с $f$, а выходной сигнал равен новому состоянию клетки.

\vspace{-30pt}
{\Large \textbf{\subsection{Окрестность фон Неймана}}} \par

Окрестность фон Неймана клетки $(i,j)$ - это сама клетка и четыре клетки, имеющие с ней общую сторону:
\[
 \mathcal{N}(i,j)=\{(i-1,\,j),\ (i,\,j-1),\ (i,\,j),\ (i,\,j+1),\ (i+1,\,j)\}.
\]
Состояние клетки $(i,j)$ в момент времени $t$ обозначим $s^t_{i,\,j}$. Окрестность фон Неймана изображена на Рис.\ref{fig:neumann}. 

\begin{figure}[H]
 \centering
 \renewcommand{\arraystretch}{2.2}
 \large
 \begin{tabular}{|*{3}{>{\centering\arraybackslash}m{2.6cm}|}}
 \hline
 & $s^t_{i-1,\,j}$ & \\ \hline
 $s^t_{i,\,j-1}$ & $s^t_{i,\,j}$ & $s^t_{i,\,j+1}$ \\ \hline
 & $s^t_{i+1,\,j}$ & \\ \hline
 \end{tabular}
 \caption{Окрестность фон Неймана клетки $(i,j)$}
 \label{fig:neumann}
\end{figure}


{\Large \textbf{\subsection{Тороидальные граничные условия}}} \par

Рассматриваются тороидальные граничные условия: противоположные края поля считаются соединёнными, то есть верхний и нижний края, а также левый и правый края попарно отождествляются.

Состояние клетки в строке $i$ и столбце $j$ в момент времени $t$ обозначим $s^t_{i,\,j}$, где $i,j\in\{1,\dots,n\}$, $n$ - размер поля. 

Рассмотрим случаи, затрагивающие границы поля (аргументы $f$ записаны в порядке: верх, лево, центр, право, низ):

\vspace{6pt}

\begin{itemize}[itemsep=-3pt]
 \item верхняя строка ($i=1$, $1<j<n$):
 \[ s^{t+1}_{1,\,j} = f\bigl(s^t_{n,\,j},\; s^t_{1,\,j-1},\; s^t_{1,\,j},\; s^t_{1,\,j+1},\; s^t_{2,\,j}\bigr) \]
 \item нижняя строка ($i=n$, $1<j<n$):
 \[ s^{t+1}_{n,\,j} = f\bigl(s^t_{n-1,\,j},\; s^t_{n,\,j-1},\; s^t_{n,\,j},\; s^t_{n,\,j+1},\; s^t_{1,\,j}\bigr) \]
 \item левый столбец ($j=1$, $1<i<n$):
 \[ s^{t+1}_{i,\,1} = f\bigl(s^t_{i-1,\,1},\; s^t_{i,\,n},\; s^t_{i,\,1},\; s^t_{i,\,2},\; s^t_{i+1,\,1}\bigr) \]
 \item правый столбец ($j=n$, $1<i<n$):
 \[ s^{t+1}_{i,\,n} = f\bigl(s^t_{i-1,\,n},\; s^t_{i,\,n-1},\; s^t_{i,\,n},\; s^t_{i,\,1},\; s^t_{i+1,\,n}\bigr) \]
\end{itemize}

В угловых клетках за границу поля выходят сразу два соседа:
\begin{itemize}[itemsep=-3pt]
 \item левый верхний угол ($i=1$, $j=1$):
 \[ s^{t+1}_{1,\,1} = f\bigl(s^t_{n,\,1},\; s^t_{1,\,n},\; s^t_{1,\,1},\; s^t_{1,\,2},\; s^t_{2,\,1}\bigr) \]
 \item правый верхний угол ($i=1$, $j=n$):
 \[ s^{t+1}_{1,\,n} = f\bigl(s^t_{n,\,n},\; s^t_{1,\,n-1},\; s^t_{1,\,n},\; s^t_{1,\,1},\; s^t_{2,\,n}\bigr) \]
 \item левый нижний угол ($i=n$, $j=1$):
 \[ s^{t+1}_{n,\,1} = f\bigl(s^t_{n-1,\,1},\; s^t_{n,\,n},\; s^t_{n,\,1},\; s^t_{n,\,2},\; s^t_{1,\,1}\bigr) \]
 \item правый нижний угол ($i=n$, $j=n$):
 \[ s^{t+1}_{n,\,n} = f\bigl(s^t_{n-1,\,n},\; s^t_{n,\,n-1},\; s^t_{n,\,n},\; s^t_{n,\,1},\; s^t_{1,\,n}\bigr) \]
\end{itemize}

{\Large \textbf{\subsection{Правило перехода}}} \par

Состояние клетки на шаге $t+1$ вычисляется по состояниям клеток окрестности на шаге $t$:
\[
 s^{t+1}_{i,\,j} = f\bigl(s^t_{i-1,\,j},\, s^t_{i,\,j-1},\, s^t_{i,\,j},\, s^t_{i,\,j+1},\, s^t_{i+1,\,j}\bigr).
\]

Окрестность из пяти клеток дает $2^5=32$ возможных комбинаций клеток. Правило перехода однозначно кодируется 32-битным числом $N$. Его двоичная запись, дополненная нулями слева до 32 разрядов, является столбцом значений булевой функции $f$ в таблице истинности.

Номер варианта $N=50045688$ в двоичной записи равен 
\[
 N = 00000010111110111010001011111000_2,
\]
Таблица истинности приведена в таблице \ref{tab:truth}.

\begin{table}[H]
\centering
\renewcommand{\arraystretch}{0.95}
\begin{tabular}{|c|c|c|c|c|c|c|}
\hline
$i$ & $S_T$ & $S_L$ & $S_C$ & $S_R$ & $S_B$ & $f$ \\
\hline
0 & 0 & 0 & 0 & 0 & 0 & 0 \\
1 & 0 & 0 & 0 & 0 & 1 & 0 \\
2 & 0 & 0 & 0 & 1 & 0 & 0 \\
3 & 0 & 0 & 0 & 1 & 1 & 0 \\
4 & 0 & 0 & 1 & 0 & 0 & 0 \\
5 & 0 & 0 & 1 & 0 & 1 & 0 \\
6 & 0 & 0 & 1 & 1 & 0 & 1 \\
7 & 0 & 0 & 1 & 1 & 1 & 0 \\
8 & 0 & 1 & 0 & 0 & 0 & 1 \\
9 & 0 & 1 & 0 & 0 & 1 & 1 \\
10 & 0 & 1 & 0 & 1 & 0 & 1 \\
11 & 0 & 1 & 0 & 1 & 1 & 1 \\
12 & 0 & 1 & 1 & 0 & 0 & 1 \\
13 & 0 & 1 & 1 & 0 & 1 & 0 \\
14 & 0 & 1 & 1 & 1 & 0 & 1 \\
15 & 0 & 1 & 1 & 1 & 1 & 1 \\
16 & 1 & 0 & 0 & 0 & 0 & 1 \\
17 & 1 & 0 & 0 & 0 & 1 & 0 \\
18 & 1 & 0 & 0 & 1 & 0 & 1 \\
19 & 1 & 0 & 0 & 1 & 1 & 0 \\
20 & 1 & 0 & 1 & 0 & 0 & 0 \\
21 & 1 & 0 & 1 & 0 & 1 & 0 \\
22 & 1 & 0 & 1 & 1 & 0 & 1 \\
23 & 1 & 0 & 1 & 1 & 1 & 0 \\
24 & 1 & 1 & 0 & 0 & 0 & 1 \\
25 & 1 & 1 & 0 & 0 & 1 & 1 \\
26 & 1 & 1 & 0 & 1 & 0 & 1 \\
27 & 1 & 1 & 0 & 1 & 1 & 1 \\
28 & 1 & 1 & 1 & 0 & 0 & 1 \\
29 & 1 & 1 & 1 & 0 & 1 & 0 \\
30 & 1 & 1 & 1 & 1 & 0 & 0 \\
31 & 1 & 1 & 1 & 1 & 1 & 0 \\
\hline
\end{tabular}
\caption{Таблица истинности функции $f$ для $N=50045688$}
\label{tab:truth}
\end{table}

\vspace{-20pt}

{\Large \textbf{\subsection{Классы Вольфрама}}} \par

Классы Вольфрама - классификация клеточных автоматов, основанная на их поведении. Вольфрам выделил 4 класса:
\begin{itemize}[itemsep=-3pt]
 \item Класс I. После некоторого количества шагов система стабилизируется: все клетки поля переходят в одинаковое состояние.
 \item Класс II. Состояния могут быть различными, однако клетки образуют наборы статичных или периодических простых структур.
 \item Класс III. Образуются сложные структуры, характер взаимодействия которых во многих случаях выглядит хаотическим.
 \item Класс IV. Сложные структуры; очень сложное, порой непредсказуемое (хаотичное) взаимодействие.
\end{itemize}

\newpage
